This paper isolates and names a form of injustice that wrongs neither what a subject \emph{has}, nor how she is \emph{recognised}, nor whether she can be \emph{heard}, but her \emph{being}---her standing as a generative being who continually constitutes herself. When intimate life is understood as a field of co-generation, in which value resides in neither party but is generated \emph{between} them, a distinctive harm becomes possible: one subject, generating too fast and too far ahead, may come to occupy the shared generative space so completely that the other loses the capacity to generate at all. The result is not that she has less but that she ceases, within the relation, to \emph{be}. The paper calls this \emb{generative effacement} and argues that it is an \emph{ontological} injustice, irreducible to the distributive, recognitive, or constitutive wrongs with which it might be confused. Its defining and most troubling feature is that it arrives not through domination but through \emph{devotion}: there is no appropriation, no coercion, no failure to hear, and the effaced party may sincerely endorse and even cherish the one who effaces her.

The harm is given an exact formal criterion. Modelling two intimate subjects as coupled phase oscillators, the paper shows that effacement is a genuine \emph{phase transition}: the weaker pole loses its own rhythm precisely when the coupling strength exceeds a threshold fixed by the difference of the two intrinsic frequencies, while the \emph{depth} of effacement, once that threshold is crossed, is set by the asymmetry of force between them. The two structures are independent, and their separation yields the central result---that even under perfectly \emph{equal} force the weaker pole is effaced by half once the threshold is crossed. Equal power does not protect it; strong coupling between unlike rhythms is sufficient to silence the slower. The good of intimacy is shown to lie not in the strength of coupling but in whether it leaves each subject turning at its own frequency.

From this follows a justice of counter-intuitive form. Because the threshold depends on the coupling and not on the stronger subject's speed, the remedy is not for the stronger to generate \emph{less}---deceleration only re-imposes its pace---but to \emph{decouple}: to release the traction by which it overwrites the other's rhythm. The paper develops this as \emb{generative reticence}, a negative and structural obligation to refrain from determining the domain in which the other generates, specifies its practice as the relinquishment of three determining channels (of the shared goal, of meaning, and of the very problems on which the relation works), and grounds it ethically in an account that extends the duty to treat another as an end from rational to generative agency, while forbidding the ethic its own paternalistic application. Justice between two generative beings, it concludes, requires not the transfer of value but the non-foreclosure of the other's capacity to generate.
